% =====================================================================
% R1.6 — "Over the seizure-onset zone, β and α diverge" — (bridge to R3)
% Closing paragraph-block of Result 1; band-dependent disease-core contrast → clinical bridge.
%
% ESTIMATOR = rho_sym (MIGRATED 2026-07-06, audit_156; same run as R1.4).
% FRAMEWORK PURITY: cophenetic objects only; RAW is baseline. α recruitment is a COPHENETIC
%   result (raw α epi_epi not concentrated → consistent with α being cophenet-only).
% Numbers (data/audit/epi_wm_stratified_rhosym/pairclass_rhosym_cohort.csv, verified 2026-07-06):
%   β epi↔epi: p_pair=0.93 (trending depleted) + matched-strength n.s. (p=0.31) → β SPARES the
%     core (directional, from R1.4).
%   α epi↔epi: obs +0.410, pair-class p_pair<0.001 CONCENTRATED + matched-strength Wilcoxon
%     p=0.005 (STRENGTHENS) → α RECRUITS the core (genuine, BOTH nulls agree). (rho_split was
%     p_pair<0.001, MS p=0.003 — like-for-like.) epi cross α also concentrated (p_pair=0.016,
%     MS p=0.003) corroborates.
% FRAMING: use R1-established descriptors (β=focal/OFC-anchored, α=diffuse/no-address) — do NOT
%   call them "memory/inference bands" here (that decomposition is R2; R1.6 precedes it). This
%   is the FIRST bridge to R3 (same organization → epileptogenic read-out).
% Guardrails: LRG framework; matched-strength named; no behavioural claim.
% Prose unwrapped (one line per paragraph); American English.
% =====================================================================

\paragraph{Over the seizure-onset zone, \(\beta\) and \(\alpha\) diverge.}
The \(\beta\) trace avoids the diseased core (above), but how the persistence treats that tissue depends on the rhythm, and the contrast is the first link between the cognitive trace and the pathology. Where \(\beta\) --- the focal, cortically anchored band --- is carried by healthy coupling and leaves the \gls{soz} out, \(\alpha\) --- the diffuse band with no cortical address --- does the reverse: the same per-pair persistence \emph{concentrates} in coupling internal to the \gls{soz} (pair-class control \(p < 0.001\)) and \emph{strengthens} rather than weakens under the strength-matched surrogate (\(p = 0.005\)), a genuine recruitment of the diseased core rather than a strength artifact. So the band that anchors in \gls{ofc} steers around the epileptogenic tissue while the band with no cortical home pulls it in (\FigRef{fig:trace2}c). The cognitive trace and the epileptogenic network are therefore not independent --- the bridge to the final result, where the same network organization that carries the cognitive trace is turned around to read out the epileptogenic tissue itself.
